\section{希格斯支与超凯勒流形}\label{sec:higgsbranch}
到目前为止，我们只考虑了超对称真空模空间中这样一个分支：矢量多重态中的标量 $\Phi$ 非零，而所有超多重态都为零。现在换一个分支，令 $\Phi=0$，但超多重态标量非零。这个分支称为希格斯支。

\subsection{希格斯支拉格朗日量的一般结构}\label{sec:hypergeneral}\index{HiggsBranch@希格斯支|hyperpage}\index{HyperkahlerManifold@超凯勒流形|hyperpage}\index{KahlerManifold@凯勒流形|hyperpage}\index{Quaternion@四元数|hyperpage}\index{HopfFibration@Hopf 纤维化|hyperpage}
首先回忆只含标量和费米子的一般 $\cN{=}1$ 理论。这样的理论可以用拉格朗日量
\begin{equation}
\int d^4\theta K(\bar\Phi_{\bar j},\Phi_i)=g_{i\bar j}\partial_\mu\phi^i\partial_\nu\bar\phi^{\bar j}+\cdots
\end{equation}
描述，其中
\begin{equation}
g_{i\bar j}=\frac{\partial^2K}{\partial\phi_i\partial\bar\phi_{\bar j}}.
\end{equation}
这定义了一个凯勒流形。特别地，该流形自然是一个复流形。这一事实在我们的形式体系中几乎是不言自明的，因为按定义，手征多重态是复值的。回忆其原因很有启发：基本超对称变换为
\begin{equation}
\delta_\alpha\phi=\psi_\alpha,\quad
\delta^\dagger_{\dot\alpha}\psi_\alpha=i\sigma^\mu_{\alpha\dot\alpha}\partial_\mu\phi
\end{equation}
一个不依赖于约定的事实是，$\delta_\alpha\delta_{\dot\alpha}$ 作用于复标量时，涉及乘以 $i$ 的运算。用 $\phi$ 的实部和虚部，可以将这一事实示意性地写成
\begin{equation}
\delta^\dagger_{\dot\alpha}\delta_\alpha
\begin{pmatrix}\Re\phi\\\
=\sigma^\mu_{\dot\alpha\alpha}\partial_\mu I
\begin{pmatrix}\Re\phi\\\
\end{equation}
其中矩阵
\begin{equation}
I=\begin{pmatrix}0&1\\-1&0\end{pmatrix}
\end{equation}
具有性质 $I^2=-1$。这就是定义 $\cN{=}1$ 理论标量流形复结构的关键矩阵。

现在考虑只由标量和费米子组成的 $\cN{=}2$ 理论。注意，这意味着不存在 $\cN{=}2$ 矢量多重态。这一理论具有两组 $\cN{=}1$ 超对称性 $\delta^{i=1,2}_\alpha$。此外，当 $|c_1|^2+|c_2|^2=1$ 时，
\begin{equation}
\delta^{(c)}_\alpha:=c_1\delta^1_\alpha+c_2\delta^2_\alpha
\end{equation}
也生成一个 $\cN{=}1$ 子超对称性。将上一段的论证应用于这个 $\cN{=}1$ 子代数，就得到矩阵
\begin{equation}
I^{(c)}=I_an^a,\qquad n^a=(\bar c_1,\bar c_2)\sigma^a\begin{pmatrix}c_1\\
\end{equation}
它们总是满足
\begin{equation}
(I^{(c)}){}^2=-1.\label{Icsq}
\end{equation}
注意，$n^a$ 为实数，且 $|n^1|^2+|n^2|^2+|n^3|^2=1$，即它们位于 $S^2$ 上。为简洁起见，记 $(I,J,K):=(I_1,I_2,I_3)$，并展开式 \eqref{Icsq}，得到关系
\begin{equation}
I^2=J^2=K^2=-1,\qquad
IJ=K=-JI,\ 
JK=I=-KJ,\ 
KI=J=-IK.\label{qcom}
\end{equation}
$I$、$J$、$K$ 的这一对易关系就是四元数的关系。切空间上具有四元数代数作用的流形称为超凯勒流形。因此，我们发现，没有无质量矢量多重态的 $\cN{=}2$ 理论的标量流形是超凯勒的。

注意，$\SU(2)_R$ 对称性作用于二重态 $(c_1,c_2)$，而它被限制在三维球面 $|c_1|^2+|c_2|^2=1$ 上。从 $(c_1,c_2)$ 到 $n^a$ 的映射 \eqref{cn} 是标准的 Hopf 纤维化 $S^3\to S^2$，指标 $a$ 按 $\SU(2)_R$ 的三重态变换。

结合第 \ref{sec:lel} 小节的分析，可以看到，一般低能 $\cN{=}2$ 理论的作用量具有形式
\begin{equation}
\int d^2\theta\frac{-\ii}{8\pi}\tau^{ij}W_{\alpha i}W^\alpha_j+cc.
+\int d^4\theta K_v(\bar a_{\bar j},a_i)
+\int d^4\theta K_h(\bar q_{\bar t},q_s)
\end{equation}
其中 $K_h(\bar q_{\bar t},q_s)$ 给出一个超凯勒流形，并且存在一个预势 $F(a_i)$，通过标准公式 \eqref{pretau}、\eqref{adi} 和 \eqref{Keq} 给出 $\tau^{ij}$ 与 $K_v$。

注意，超多重态部分与矢量多重态部分完全退耦。对紫外规范耦合的依赖隐含在矢量多重态部分之中。这意味着超多重态部分不能受到依赖规范耦合的量子修正。

\subsection{再论超多重态}\label{sec:hyperrevisited}\index{Hypermultiplet@超多重态|hyperpage}\index{HalfHypermultiplet@半超多重态|hyperpage}\index{PseudoRealRepresentation@赝实表示|hyperpage}\index{FlavorSymmetry@味对称性|hyperpage}\index{Bifundamental@双基本表示|hyperpage}
从这里的观点重新考察第 \ref{sec:hyperdef} 小节引入的完整超多重态与半超多重态的结构。首先回忆紧群不可约表示的类型：
\begin{equation}
\begin{array}{rll}
\text{复表示}&\text{若 $R\not\simeq\bar R$；}&\\
\text{实表示}&\text{若 $R\simeq\bar R$：}&\begin{cases}
\text{严格实表示}&\text{若不变张量 $\delta_{ij}$ 对称，}\\
\text{赝实表示}&\text{若不变张量 $\epsilon_{ij}$ 反对称。}
\end{cases}
\end{array}
\end{equation}

在含有若干实标量 $\phi^i$（$i=1,\ldots,n$）的非超对称理论中，若存在表示 $G$ 的李代数的实 $n\times n$ 矩阵 $T^a$（$a=1,\ldots,\dim G$），就可以有味对称性 $F$ 或规范对称群 $G$ 对这些标量的作用：
\begin{equation}
F,G\curvearrowright\bR^n.
\end{equation}
这一表示显然具有不变的对称张量 $\delta_{ij}$，因为它作用于具有动能项 $\delta_{ij}\partial_\mu\phi^i\partial_\mu\phi^j$ 的 $n$ 个实标量。因而，此表示是严格实的。

在含有若干实标量 $\phi^i$（$i=1,\ldots,n$）以及 $n/2$ 个 Weyl 费米子的 $\cN{=}1$ 超对称理论中，超对称性要求存在满足 $I^2=-1$ 的矩阵 $I$。味对称性 $F$ 和规范对称性 $G$ 的作用必须与此矩阵 $I$ 对易：
\begin{equation}
F,G\curvearrowright\bR^n\curvearrowleft I.
\end{equation}
可以规定，复数 $a+bi$ 通过矩阵 $a+bI$ 作用于实标量。于是，标量空间变为一个复向量空间，而对称性的作用保持复结构。因此，存在 $m=n/2$ 个手征多重态 $\Phi^s$（$s=1,\ldots,m$），$F$ 和 $G$ 都通过表示其李代数的 $m\times m$ 复矩阵作用于它们。可以将这一情形概括为
\begin{equation}
F,G\curvearrowright\bC^m.
\end{equation}

在含有若干实标量 $\phi^i$（$i=1,\ldots,n$）以及 $n/2$ 个 Weyl 费米子的 $\cN{=}2$ 超对称理论中，超对称性要求存在满足对易关系 \eqref{qcom} 的矩阵 $I$、$J$、$K$。味对称性 $F$ 和规范对称性 $G$ 的作用必须与 $I$、$J$、$K$ 对易：
\begin{equation}
F,G\curvearrowright\bR^n\curvearrowleft I,J,K.
\end{equation}
可以规定，四元数 $a+bi+cj+dk$ 通过矩阵 $a+bI+cJ+dK$ 作用于实标量。于是，标量空间变为一个四元数向量空间，而对称性的作用保持四元数结构。这自动要求 $n$ 为四的倍数，即 $n=4\ell$。$F$ 和 $G$ 都通过表示其李代数的 $\ell\times\ell$ 四元数矩阵作用于它们。总结起来，
\begin{equation}
F,G\curvearrowright\bH^\ell
\end{equation}
其中 $\bH$ 是四元数除环。

由于四元数在物理学家中并不十分常用，我们通常只使用 $a+bI$，将实标量视为复标量。此时得到一个维数为 $2\ell$ 的复向量空间，其中有 $2\ell$ 个复标量 $\Phi^s$（$s=1,\ldots,2\ell$），味对称性 $F$ 和规范对称性 $G$ 通过复表示 $\tilde R$ 作用于它们。矩阵 $J+iK$ 随后确定一个 $2\ell\times2\ell$ 反对称矩阵 $\epsilon_{st}$，它在 $F$ 和 $G$ 的作用下不变。这意味着 $\tilde R$ 是一个赝实表示：
\begin{equation}
F,G\curvearrowright\bC^{2\ell}\curvearrowleft\epsilon_{st}
\end{equation}
这就是第 \ref{sec:hyperdef} 小节末尾简要引入的表示 $\tilde R$ 中的半超多重态。

从这一观点看，半超多重态比完整超多重态更基本，后者可以如下给出。取 $F\times G$ 的任意一个 $m$ 维复表示 $R$，以 $i=1,\ldots,m$ 为其指标。存在一个不变张量 $\delta_{i\bar j}$。令 $\tilde R=R\oplus\bar R$。其指标为 $s=1,\ldots,n,\bar1\ldots,\bar n$，并自动具有一个反对称不变张量
\begin{equation}
\epsilon_{st},\qquad\epsilon_{ij}=\epsilon_{\bar i\bar j}=0,\ \epsilon_{i\bar j}=\delta_{i\bar j}=-\epsilon_{\bar ji}.
\end{equation}
基于表示 $\tilde R$ 的半超多重态，就是表示 $R$ 中的完整超多重态。

% [Tutor Note: 本段前文取表示 $R$ 的维数为 $m$、其指标为 $i=1,\ldots,m$，而这里给 $\tilde R$ 的指标却写作 $s=1,\ldots,n,\bar1\ldots,\bar n$，与 $m$ 不符，疑为笔误（应为 $m$）；按底本照原文排。]

具体而言，考虑四个实标量。这个系统具有自然的对称性 $\SO(4)\simeq\SU(2)_l\times\SU(2)_r$。再加入两个 Weyl 费米子，它们具有自然的对称性 $\SU(2)_l$。于是，整个系统具有 $\cN{=}2$ 超对称性，其中 $\cN{=}2$ 代数的 $\SU(2)_R$ 对称性，就是作用于标量的 $\SU(2)_r$。对称性 $\SU(2)_l$ 可以作为味对称性，也可以作为规范对称性。整个系统恰好由一个完整超多重态组成，或等价地，由一个 $\SU(2)_l$ 二重态中的半超多重态组成。

接下来，令 $i=1,\ldots,n$ 和 $a=1,\ldots,m$ 分别为 $\U(n)$ 和 $\U(m)$ 对称性的指标。$\cN{=}1$ 手征多重态 $Q_{i\bar a}$、$\tilde Q_{\bar ia}$ 随即构成 $\U(n)\times\U(m)$ 双基本表示中的一个 $\cN{=}2$ 超多重态。当 $\U(n)$ 被视为规范对称性时，$\U(m)$ 就成为味对称性。

另一个典型构造，是取 $i=1,\ldots,2n$ 为 $\Sp(n)$ 对称性的指标，取 $a=1,\ldots,m$ 为 $\SO(m)$ 对称性的指标。考虑 $\cN{=}1$ 手征多重态 $Q_{ia}$。将指标对 $ia$ 视为单个指标 $s=(ia)$，取值从 $1$ 到 $2nm$。此系统具有反对称不变张量 $\epsilon_{st}=\epsilon_{(ia)(jb)}=J_{ij}\delta_{ab}$，因此它们构成一个具有 $\Sp(n)\times\SO(m)$ 对称性的超多重态，而这一对称性与超代数对易。当 $\Sp(n)$ 成为规范对称性时，$\SO(m)$ 就成为味对称性，反之亦然。这就解释了如下事实：若有 $n$ 个超多重态处于规范群 $\SO(m)$ 的矢量表示中，就具有 $\Sp(n)$ 味对称性；若有 $m$ 个半超多重态处于规范群 $\Sp(n)$ 的基本表示中，就具有 $\SO(m)$ 味对称性。

\subsection{超凯勒商}\label{sec:hkq}\index{HyperkahlerQuotient@超凯勒商|hyperpage}\index{ComplexifiedGaugeGroup@复化规范群|hyperpage}
回到希格斯支的研究。对含有 $N_f$ 个味的 $\SU(N)$ 规范理论，定义希格斯支的方程已在第 \ref{sec:classicalvac} 小节给出，见式 \eqref{Feq} 和 \eqref{tripleteq}。现在将其写成一般情形的形式。

考虑规范群为 $G$ 的 $\cN{=}2$ 规范理论，并有一个处于表示 $R$ 中的超多重态 $(Q^i,\tilde Q_i)$。这里，指标 $i=1,\ldots,\dim R$ 用于超多重态，而指标 $a=1,\ldots,\dim G$ 用于伴随表示。希格斯支由下式给出：
\begin{equation}
\left\{
\begin{array}{r@{=0}}
(Q^iQ^\dagger{}_j-\tilde Q^\dagger{}^i\tilde Q_j)T^a{}_i^j\\
\Re Q^i\tilde Q_jT^a{}_i^j\\
\Im Q^i\tilde Q_jT^a{}_i^j
\end{array}
\right\}\Bigm/(\text{按规范群作用作等同})\label{hkq}
\end{equation}
其中 $T^a{}_i^j$ 是 $G$ 的代数在表示 $R$ 中的矩阵。

在希格斯支的一般点上，没有剩余的无质量矢量多重态。根据前面的一般分析，我们知道它们构成一个超凯勒流形。在数学和物理文献中，构造 \eqref{hkq} 称为超凯勒商构造，最初由文献 \cite{Hitchin:1986ea} 发现。任意超凯勒流形的实维数总是四的倍数。让我们在当前情形中检验这一点。假设最初的超多重态由 $4m$ 个实标量组成。对每一个 $I$、$J$、$K$，D 项条件都施加 $\dim G$ 个实约束。随后，再按 $G$ 的作用作等同。因此，取商后的实维数为
\begin{equation}
4m-3\dim G-\dim G=4(m-\dim G)\label{realdim}
\end{equation}

如果只关心全纯结构，可以舍去 D 项方程，转而按复化规范群作等同：
\begin{equation}
\left\{Q^i\tilde Q_jT^a{}_i^j=0\right\}\Bigm/(\text{按\emph{复化}规范群作用作等同}).\label{hsq}
\end{equation}
注意，在 $\cN{=}1$ 超场形式中，如果不将矢量超场置于 Wess-Zumino 规范，这是一种更自然的形式。证明式 \eqref{hsq} 与 \eqref{hkq} 相等的基本思想，是在复化规范群的每一条轨道中将 $|D|^2$ 最小化。最小化条件给出 $D=0$，从而恢复式 \eqref{hkq}。这一粗略分析也表明，更准确地说，需要从式 \eqref{hsq} 中去除所谓的不稳定轨道；在这些轨道上，不存在使 $|D|^2$ 取最小值的点。

在这种方法中，从 $2m$ 个复标量出发，施加 $\dim G$ 个复约束，再按复化规范群 $G_\bC$ 的作用作等同，去掉 $\dim G$ 个复维度。最终，商空间中有
\begin{equation}
2m-\dim G-\dim G=2(m-\dim G)\label{complexdim}
\end{equation}
个复标量。这与刚刚在式 \eqref{realdim} 中得到的结果相容。若计算四元数维数，则公式就是
\begin{equation}
m-\dim G.
\end{equation}

\subsubsection{含有一个带荷超多重态的 $\U(1)$ 规范理论}
考虑两个例子。首先取一个 $\cN{=}2$ $\U(1)$ 规范理论，其中有两个荷为 $\pm1$ 的超多重态 $(Q_i,\tilde Q^i)$，这里 $i=1,2$。在上述表达式中，$m=2$、$\dim G=1$，因此预期得到一个复二维的希格斯支。首先显式确定希格斯支。F 项方程为
\begin{equation}
Q_1\tilde Q^1+Q_2\tilde Q^2=0.
\end{equation}
于是有
\begin{equation}
(Q_1,\tilde Q^1)=(z,\tilde zt),\quad
(Q_2,\tilde Q^2)=(\tilde z,-zt)
\end{equation}
其中 $z$、$\tilde z$ 和 $t$ 为某些复数。D 项方程 $|Q_1|^2+|Q_2|^2-|\tilde Q^1|^2-|\tilde Q_2|=0$ 给出
\begin{equation}
|z|^2+|\tilde z|^2=|t|^2(|z|^2+|\tilde z|^2).
\end{equation}

% [Tutor Note: 原文此处 D 项方程印作 $|Q_1|^2+|Q_2|^2-|\tilde Q^1|^2-|\tilde Q_2|=0$，最后一项脱去平方（v3 已补作 $|\tilde Q_2|^2$），与其余三项不一致；按底本 v2 照原文排。]
因此有 $|t|=1$。可以利用 $\U(1)$ 规范旋转，通过要求
\begin{equation}
\Arg z=\Arg(\tilde zt).
\end{equation}
来几乎完全消去 $t$。这仍未固定将 $Q_i$、$\tilde Q^i$ 乘以 $-1$ 的 $\U(1)$ 规范变换，它将 $(z,\tilde z)$ 映为 $(-z,-\tilde z)$。因此，希格斯支为
\begin{equation}
\bC^2/\bZ_2=\{(z,\tilde z)\in\bC^2\}/(z,\tilde z)\leftrightarrow(-z,-\tilde z).\label{c2z2}
\end{equation}
图 \ref{fig:c2z2} 给出了一个不太精确的示意描述。
\begin{figure}[h]
\[
\inc{c2z2}
\]
\caption{$\bC^2/\bZ_2$ 的一个不太精确的示意图。\label{fig:c2z2}}
\end{figure}

利用复描述 \eqref{hsq}，可以用另一种方式得到同一个希格斯支。不必将复化规范群连接起来的点作等同，而可以直接考虑在其作用下不变的坐标组合。在这里，$Q_i$ 的荷为 $+1$，$\tilde Q^i$ 的荷为 $-1$。因此，规范不变量是 $Q_i\tilde Q^j$，其中 $i$、$j$ 可以任意选取。还需要施加
\begin{equation}
Q_1\tilde Q^1+Q_2\tilde Q^2=0,
\end{equation}
这一条件。总共得到三个组合
\begin{equation}
A=Q_1\tilde Q^2,\quad
B=Q_2\tilde Q^1,\quad
C=iQ_1\tilde Q^1=-iQ_2\tilde Q^2.
\end{equation}
它们满足一个显然的关系
\begin{equation}
AB=C^2.\label{abc}
\end{equation}
三个变量 $A$、$B$、$C$ 加上上述一个关系，给出一个复二维空间。这就是希格斯支。

从式 \eqref{c2z2} 中 $\bC^2/\bZ_2$ 的定义出发，也可以得到这一描述。在 $(z,\tilde z)\leftrightarrow(-z,-\tilde z)$ 下不变的 $z$、$\tilde z$ 的组合为
\begin{equation}
A=z^2,\quad B=\tilde z^2,\quad C=z\tilde z
\end{equation}
它们满足同一个关系 \eqref{abc}。因此，作为复流形，它们是同一个空间。

\subsubsection{含有两个二重态超多重态的 $\SU(2)$ 规范理论}
作为第二个例子，考虑具有 $N_f$ 个二重态表示完整超多重态的 $\cN{=}2$ $\SU(2)$ 规范理论。用 $\cN{=}1$ 手征多重态表示，有
\begin{equation}
Q_i^a,\quad\tilde Q^i_a\qquad(a=1,2;\quad i=1,\ldots,N_f)
\end{equation}
由于对 $\SU(2)$ 规范理论而言，二重态与反二重态表示相同，也可以将其记为
\begin{equation}
Q_I^a,\qquad(a=1,2;\quad I=1,\ldots,2N_f)
\end{equation}
这使 $\SO(2N_f)$ 味对称性更加明显。

共有 $4N_f$ 个复标量，且 $\dim\SU(2)=3$。因而，希格斯支的复维数为
\begin{equation}
4N_f-2\cdot3.\label{dimform}
\end{equation}
所以 $N_f=1$ 时没有希格斯支，而 $N_f=2,3,4$ 时，预期分别有复维数为 $2$、$6$、$10$ 的希格斯支。下面更详细地研究 $N_f=2$ 的情形。

$Q_I^a$ 的规范不变组合为
\begin{equation}
M_{IJ}=Q_I^aQ_J^b\epsilon_{ab}.\label{su2mesons}
\end{equation}
左端在交换 $I$ 与 $J$ 时自动反对称。F 项方程为
\begin{equation}
Q_I^{(a}Q_J^{b)}\delta^{IJ}=0.\label{su2Dterm}
\end{equation}

对于 $\SO(2N_f)=\SO(4)$，可以将 $\SO(4)$ 的反对称矩阵 $M_{IJ}$ 分解为自对偶部分和反自对偶部分。等价地，利用 $\SO(4)\simeq\SU(2)_u\times\SU(2)_v$，$M_{IJ}$ 分解为 $\SU(2)_u$ 的三重态 $M_{(\alpha\beta)}$ 和 $\SU(2)_v$ 的三重态 $M_{(\dot\alpha\dot\beta)}$，其中 $\alpha,\beta=1,2$ 与 $\dot\alpha,\dot\beta=1,2$ 分别是 $\SU(2)_{u,v}$ 的二重态指标。指标 $I$ 本身可以视为一对指标：$I=(\alpha\dot\alpha)$。于是，所讨论的超多重态可以写成
\begin{equation}
Q_{a\alpha\dot\alpha},\qquad a,\alpha,\dot\alpha=1,2
\end{equation}
从而使 $\SU(2)^3$ 对称性显现出来。此时
\begin{align}
M_{\alpha\beta}&=Q_{a\alpha\dot\alpha}Q_{b\beta\dot\beta}\epsilon^{ab}\epsilon^{\dot\alpha\dot\beta},\\
M_{\dot\alpha\dot\beta}&=Q_{a\alpha\dot\alpha}Q_{b\beta\dot\beta}\epsilon^{ab}\epsilon^{\alpha\beta}
\end{align}
分别是 $M_{IJ}$ 的自对偶和反自对偶部分。F 项方程可以写为
\begin{equation}
Q_{a\alpha\dot\alpha}Q_{b\beta\dot\beta}\epsilon^{\alpha\beta}\epsilon^{\dot\alpha\dot\beta}=0.
\end{equation}

利用这一描述，不难检验
\begin{align}
M_{\alpha\beta}M_{\gamma\delta}\epsilon^{\alpha\gamma}\epsilon^{\beta\delta}&=0,\label{c2z2ABC}\\
M_{\dot\alpha\dot\beta}M_{\dot\gamma\dot\delta}\epsilon^{\dot\alpha\dot\gamma}\epsilon^{\dot\beta\dot\delta}&=0,\label{c2z2XYZ}\\
M_{\alpha\beta}M_{\dot\alpha\dot\beta}&=0.\label{ABCperpXYZ}
\end{align}
通过定义
\begin{align}
A&=M_{11},&B&=M_{22},&C&=M_{12}=M_{21};\\
X&=M_{\dot1\dot1},&Y&=M_{\dot2\dot2},&Z&=M_{\dot1\dot2}=M_{\dot2\dot1}.
\end{align}
结构会更加清楚。关系 \eqref{c2z2ABC} 与 \eqref{c2z2XYZ} 给出
\begin{equation}
AB=C^2,\qquad XY=Z^2,
\end{equation}
而关系 \eqref{ABCperpXYZ} 意味着两个向量 $(A,B,C)$ 和 $(X,Y,Z)$ 不能同时非零。
\begin{figure}[h]
\[
\inc{c2z2c2z2}
\]
\caption{$\bC^2/\bZ_2\wedge\bC^2/\bZ_2$ 的一个不太精确的示意图。\label{fig:c2z2copies}}
\end{figure}

因此可见，希格斯支具有图 \ref{fig:c2z2copies} 所示的结构：存在两个 $\bC^2/\bZ_2$ 副本，分别由变量 $A,B,C$ 和 $X,Y,Z$ 描述。当真空位于由一组变量 $(A,B,C)$ 描述的某个 $\bC^2/\bZ_2$ 上时，另一组变量被迫为零，反之亦然。因此可以说，两个 $\bC^2/\bZ_2$ 副本共用一个原点，在那里 $A,B,C$ 和 $X,Y,Z$ 全部为零。正如预期，希格斯支的复维数为二。

回忆一下，我们将味对称性 $\SO(4)$ 分解为 $\SU(2)_u\times\SU(2)_v$。向量 $(A,B,C)$ 和 $(X,Y,Z)$ 分别是 $\SU(2)_u$ 与 $\SU(2)_v$ 的三重态。因此，$\OO(4)\supset\SO(4)$ 中的味宇称交换组成希格斯支的两个 $\bC^2/\bZ_2$ 副本。
